\chapter{Introduction to the Finite Element Method}
\label{ch:intro}

\section{Purpose of the Finite Element Method}

Many engineering problems are governed by differential equations together with
boundary conditions. For simple geometries and idealized loading, these
equations may have closed-form analytical solutions. In practical design,
however, the geometry, material behavior, loading, and supports are often too
complex for an exact solution.

The Finite Element Method (FEM) is a systematic numerical method for replacing
a continuous problem by a finite set of algebraic equations. Instead of solving
for an unknown field at infinitely many points, FEM solves for a finite number
of unknowns at selected points called nodes. The solution between nodes is then
interpolated by simple functions.

FEM is widely used in structural, mechanical, civil, aerospace, biomedical, and
multiphysics engineering because it combines physical modeling with a practical
computational procedure.

\section{Historical Background}

The modern finite element method developed in the 1950s from matrix structural
analysis and the need to solve complex aerospace structures. Turner, Clough,
Martin, and Topp~\cite{zienkiewicz2005fem} published an important early paper
in 1956 describing the stiffness approach for plane stress analysis. The term
\emph{finite element} was introduced by Clough in 1960.

Since then, FEM has grown into one of the central tools of computational
engineering. It is now implemented in commercial and open-source software for
linear and nonlinear statics, heat transfer, dynamics, fluid flow, and coupled
field problems.

\section{Basic Idea}

The central idea of FEM is discretization. A domain is divided into small,
simple parts called \textbf{elements}. Elements meet at \textbf{nodes}, where the
primary unknowns are stored. Within each element, the unknown field is
approximated using interpolation functions called \textbf{shape functions}.

For a structural displacement problem, the final algebraic system commonly has
the form
\begin{equation}
    \stiff \, \vect{u} = \vect{f}
    \label{eq:global_system}
\end{equation}
where $\stiff$ is the global stiffness matrix, $\vect{u}$ is the vector of
unknown nodal displacements, and $\vect{f}$ is the global load vector.

Although \cref{eq:global_system} looks compact, it represents the combined
effect of element behavior, material properties, geometry, loading, and boundary
conditions. Much of FEM is the careful construction of this system.

\section{General FEM Procedure}

A typical finite element analysis follows these steps:

\begin{enumerate}
    \item \textbf{Define the physical problem:} Identify the domain, material
    properties, loads, and boundary conditions.
    \item \textbf{Choose the element type:} Select elements appropriate for the
    physics and expected behavior of the solution.
    \item \textbf{Discretize the domain:} Divide the domain into a mesh of
    finite elements connected at nodes.
    \item \textbf{Form element equations:} Derive the element stiffness matrix
    and element load vector.
    \item \textbf{Assemble the global system:} Combine all element equations
    into the global matrix equation.
    \item \textbf{Apply boundary conditions:} Enforce prescribed values and
    include natural loading terms.
    \item \textbf{Solve the algebraic system:} Compute the unknown nodal
    quantities.
    \item \textbf{Post-process the results:} Recover derived quantities such as
    strains, stresses, reactions, or internal forces.
\end{enumerate}

\section{Role of Approximation}

FEM is an approximate method. The accuracy of the result depends on the element
formulation, mesh density, interpolation order, numerical integration, and the
quality of the physical model. A refined mesh or higher-order interpolation can
improve accuracy, but it also increases computational cost.

Good finite element modeling therefore requires engineering judgment. The goal
is not only to obtain numbers, but to obtain numbers that are meaningful for the
problem being studied.

\section{Organization of This Book}

This book begins with weighted residual methods because they provide a clear
mathematical bridge from differential equations to finite element equations. It
then develops one-dimensional elements, truss and beam elements, and
two-dimensional continuum elements.

Each chapter emphasizes the connection between theory, element formulation, and
implementation. MATLAB examples in the appendix show how the main algorithms
can be programmed directly.

\section{Notation Used in This Book}

Throughout this book, we use the following conventions:

\begin{center}
\begin{tabular}{@{}lll@{}}
    \toprule
    Symbol & Meaning & Example \\
    \midrule
    $\vect{u}$ & Column vector (bold italic) & displacement vector \\
    $\matr{K}$ & Matrix (bold italic) & stiffness matrix \\
    $(\cdot)\transpose$ & Transpose & $\matr{K}\transpose$ \\
    $(\cdot)\inv$ & Inverse & $\matr{K}\inv$ \\
    $\dd$ & Differential operator & $\dd x$, $\dd V$ \\
    $\delta$ & Virtual quantity & $\delta u$, $\delta W$ \\
    \bottomrule
\end{tabular}
\end{center}
